\documentclass[11pt]{article}
\usepackage[margin=0.8in]{geometry}
\usepackage[T1]{fontenc}
\usepackage{amsmath,booktabs,xurl,hyperref}
\setlength{\emergencystretch}{3em}
\title{Closed-Loop Outcomes After Removing Solver-Result Admission}
\author{Pan Dynamics Collision Avoidance Research}
\date{October 1, 2026}
\begin{document}
\maketitle
\section{What was established}
The prior saved-frame replay established that eleven formerly rejected CLF
results could be returned, not that eleven complete scenarios succeeded.
This experiment reruns all fourteen scenarios from their initial states with
controller commit \path{0466fd95800e80d53434cf90b8021da5989932b0}.
The controller was unchanged during the campaign.

Of the eleven former CLF-failure cases, 0 confirmed nominal recovery and 0 completed collision-free avoidance plus recovery. Across all fourteen cases, 2 showed strict geometric overlap, 6 stopped with a controller failure, and 8 reached the 80-s observation cap without recovery. Completing a duration is not completing recovery.


\section{Conditions and outcomes}
Seven existing threatening scenarios were run at 8 and 15 m/s. Each starts
separated and collides under ideal given-path constant-speed cruise. Inputs
are exact state and fixed known target motion, without observation noise,
delay or road constraints. The configured collision buffer is 6 mm; sample
and prediction integration steps are 50 ms; total horizons are 68/76 holds.
No objective, initialization, trust region, terminal set or solver tolerance
was changed. There are two optimizer calls per returned hold. Two discarded
no-target warm-ups precede each speed group. No random draws are used.

Plant evolution uses ODE45 with relative/absolute tolerances $10^{-11}$ and
$10^{-12}$, with 31 geometry samples per hold. Vectorized NumPy rectangle geometry independently recomputes clearance and
strict overlap from saved states. Each case also compares 64 distributed
samples and its clearance/overlap extrema against the existing scalar Python
geometry, with agreement required within $10^{-8}$ m.
This is sampled simulation evidence, not a continuous-plant certificate.
There are no impact dynamics, so trajectories after first overlap do not
represent physical post-impact vehicle motion.

Recovery requires at least five consecutive seconds within errors
\[(0.1\text{ m},1^\circ,0.1\text{ m/s},0.05\text{ m/s},0.01\text{ rad/s})\]
for lateral position, heading, speed, lateral velocity and yaw rate, starting
no earlier than 8 s. An 80-s cap measures extended behavior; nonrecovery at
that cap is not a proof of infinite-time divergence. Large persistent or
growing errors are reported explicitly rather than called successful slow
recovery.

\begin{center}\small
\begin{tabular}{r l r r r l}
\toprule
Speed & Scenario & Last time (s) & Min gap (m) & Final lateral (m) & Outcome\\
\midrule
8 & headOn & 2.35 & 0.6152 & 4.61 & Stopped \\
8 & acceleratingHeadOn & 80.00 & 0.4682 & 354.07 & 80-s cap \\
8 & brakingLead & 11.40 & 1.1143 & -2.02 & Stopped \\
8 & crossing & 1.40 & 2.3205 & 0.01 & Stopped \\
8 & turningCrossing & 80.00 & 0.4354 & -296.94 & 80-s cap \\
8 & curvedHeadOn & 1.35 & 2.8554 & 2.46 & Stopped \\
8 & curvedCrossing & 80.00 & 0.5661 & -224.23 & 80-s cap \\
15 & headOn & 80.00 & 0.0000 & 725.12 & 80-s cap; overlap \\
15 & acceleratingHeadOn & 80.00 & 0.0000 & 727.14 & 80-s cap; overlap \\
15 & brakingLead & 23.95 & 0.8997 & -2.64 & Stopped \\
15 & crossing & 80.00 & 0.4775 & -628.85 & 80-s cap \\
15 & turningCrossing & 80.00 & 0.0471 & -728.69 & 80-s cap \\
15 & curvedHeadOn & 0.10 & 29.5064 & -0.00 & Stopped \\
15 & curvedCrossing & 80.00 & 0.0342 & -373.19 & 80-s cap \\
\bottomrule
\end{tabular}\end{center}

For 8-m/s acceleratingHeadOn, lateral error grows from 35.30 m at 10 s to
354.07 m at 80 s. The final heading error is about $0.758$ rad and the speed
error about $-0.994$ m/s. Its final CLF slack is about 109169, and even its
last affine CLF value increases over the current value. Thus the controller
is continuing with positive returned CLF slack, not exhibiting the
requested zero-slack dissipation. For 8-m/s turningCrossing and curvedCrossing,
80-s lateral errors are about $-296.94$ and $-224.23$ m. These findings do not
isolate every mechanism behind failed nonlinear recovery.

\section{Why the original three PCBF problems have no solution}
The three original PCBF-failure prefixes reproduce to numerical precision.
Removing result admission cannot produce a vector when the convex problem
itself is inconsistent. Numerical Farkas witnesses and constraint-removal
experiments are preserved in
\path{report/TWO_STAGE_FAILURE_CAUSES_20261001/}.

\subsection{8-m/s crossing and curvedHeadOn}
The next reference is a shifted affine state/input prediction. It is not
regenerated by rolling the inputs from the new measured state. Therefore
relinearization introduces dynamics defects and an initial mismatch, while
hard local boxes permit only limited corrections, including steering changes
of $\pm0.075$ rad and braking-ratio changes of $\pm0.125$ relative to the
reference. These are extra local restrictions, not actuator limits.

For both cases, deleting collision rows or terminal constraints does not
restore feasibility. Removing either the state trust bounds or input trust
bounds while retaining physical limits does. The crossing case is already
inconsistent with only dynamics and trust boxes. The curvedHeadOn conflict
also involves physical state bounds, including the positive-speed floor.
The structural defect is reuse of a dynamically inconsistent affine reference
inside a hard correction neighborhood. Perfect state observations and removal
of road boundaries do not remove that defect.

\subsection{15-m/s curvedHeadOn: premature encounter dismissal and hard geometry}
An upstream defect is visible in the unchanged range logic. Measured rectangle
clearance is 31.789 m initially, 30.648 m at 0.05 s, and 29.506 m at 0.10 s:
the target is approaching. Nevertheless, localBeyondRange classifies the
first two calls as already departed. localFlowSeed initializes its exited
flag from this current-distance test, and localFormulate sets departure to
zero, suppressing future collision rows even though the target will enter
the interaction region. The movingTargetFlow metadata label alone does not
mean its flow guidance was active.

At 0.10 s, collision constraints activate on the reused cruise reference;
its basic finite/domain reuse test does not reconstruct a target-aware seed.
Hard tail separation normals around absolute times 1.45/1.75 s point nearly
opposite longitudinal directions. These rows and the remaining dynamics and
bounds define an infeasible local passage. Removing terminal constraints or
all trust bounds does not repair it; removing hard tail collision rows does.
A current-distance threshold must not be interpreted as proof of future
conflict clearance for an approaching known target.

The previous two uniform lateral-normal probes both failed, so no claim is
made that changing one normal is sufficient. An additional diagnostic here
restarts initialization at each original failure state with no prior affine
trajectory, retaining the physical state, held input, target observation,
configuration and two-stage structure. This changes the reference/epoch
initialization, not the production algorithm. Its outcomes are:
8-m/s crossing: returned a numerical plan. 8-m/s curvedHeadOn: returned a numerical plan. 15-m/s curvedHeadOn: returned a numerical plan.
All three restarted calls returned solver flags $(1,1)$; their PCBF slack
sums were approximately $8.39\times10^{-7}$, $7.35\times10^{-7}$ and
$1.04\times10^{-6}$, respectively. Thus changing the reference construction
restores numerical solvability of these local problems at the same physical
states. It does not prove nonlinear physical feasibility.
These single-call probes do not demonstrate complete safe avoidance or
recovery and are not executable fallback policies.

\section{Runtime, scope and artifacts}
There were 13611 returned holds. Warmed full-controller median and maximum are 65.984 and 477.619 ms; the maximum including failed calls is 477.619 ms. There were 1655 returned calls and 0 failed calls over 100 ms. These times exclude offline replay and geometry processing. They are measurements, not deadline guarantees.


Online affine/nonlinear result admission remains absent. Residuals not measured
online remain null in exports; this experiment assesses physical sampled
clearance and recovery without inventing affine-feasibility claims. No
controller implementation, solver or configuration change was made here.
Unit tests were not repeated for unchanged production source.

The accompanying data bundle contains all fourteen summaries, compact frame
metrics, first-overlap details, range/seed diagnostics, executed methods,
actual console exports, and hashes of source, native kernels and raw files.
The source snapshot and large MAT/JSON replay data remain in the documented
local cache. Logs distinguish controller failures from duration completion.
\end{document}
